% 第 17 章。数字：checks/check_ch17.py；引文：data/snapshots/toku_facts.csv（清代史料逐字核对见
% data/fetch/verify_quotes.py，研究文献见 data/fetch/verify_web_quotes.py；Suzuki 的表 1—5 按页面图像读，见快照元数据）；
% 数据：gpih_iwahashi_rice、gpih_gamble_beijing（data/fetch/gpih_east_asia.py）；图：figsrc/fig_ch17.py。
\chapter{德川日本与清代中国：东亚的前现代市场}\label{ch:tokugawa}

\begin{chapterguide}
\guidequestion{东亚的前现代市场经济怎样运转？一个以米计税的国家怎样管理货币与物价？}
\begin{guidegoals}
\item 描述德川日本的石高制、金银铜三货制与元禄改铸（1695）的通胀效应；
\item 解释大阪堂岛米市（1730 年获幕府认可）中远期交易的结构；
\item 分析清代银钱二元体系，以及“银贵钱贱”的成因；
\item 比较东亚与同期欧洲在金融制度上的异同。
\end{guidegoals}
\guideprereq{\chref{ch:songming}\enspace\nameref{ch:songming}；\chref{ch:silver}\enspace\nameref{ch:silver}}
\end{chapterguide}

1695 年，德川幕府开始把市面上的金银币收回重铸。新铸的元禄小判与旧的庆长小判一样重，都是 17.85 克，面值都是 1 两（ryō），含金却从 84.291\% 降到 57.366\%，每两所含的纯金少了将近三分之一\cite{suzuki2021}。江户的米价随即上涨：按岩桥胜（Iwahashi Masaru）整理的价格，1685—1694 年平均每石 0.79 两，1696—1710 年平均 1.15 两，涨了 46\%\cite{iwahashi1981}。一个半世纪以后，清朝的曾国藩（1811—1872）在奏疏里诉说的是另一种钱的麻烦：东南产米之区，“大率石米买钱三千。自古迄今。不甚悬远。昔日两银换钱一千。则石米得银三两。今日两银换钱二千。则石米仅得银一两五钱”；朝廷收的税没有变，“而小民暗加一倍之赋”\cite{xubiansheng}。

两段话说的是同一类问题：一个国家同时用两三种钱，钱的成色或彼此的比价一变，物价、税负与债务的实际分量都跟着变。德川日本按米计算土地与赋税，却要管理金、银、铜三种货币；清代中国按银两计税，百姓日常用的却是铜钱。两国还都有发达的市场：大阪的堂岛米市在 1730 年得到幕府承认，Schaede 认为，第一个“组织完备、满足现代金融研究所列全部技术标准的期货交易所”可以追溯到 18 世纪的日本\cite{schaede1989}。

这一章回答导引框里的问题：东亚的前现代市场经济怎样运转？一个以米计税的国家怎样管理货币与物价？依次讨论：18 世纪初两国的人口、赋税与货币（\secref{sec:toku-world}）；以米计税与以银计税，各让谁承担价格的风险，东西日本的两种计价怎样抬高江户的物价（\secref{sec:toku-koku}）；元禄改铸为什么引起通胀，新井白石怎样复原币制（\secref{sec:toku-genroku}）；堂岛米市的米票与期货（\secref{sec:toku-dojima}）；清代的银钱比价怎样从“钱贵”变成“银贵”（\secref{sec:toku-qing}）；当时与今天的人怎样解释银贵（\secref{sec:toku-debate}）；东亚与同期欧洲的异同（\secref{sec:toku-europe}）。\secref{sec:toku-test} 用米价与工资的数据检验本章的两个模型。本章引用的日文史料（室鸠巢、荻生徂徕、太宰春台、新井白石）都转引自 Suzuki（2021）；外文引文的中文都是本书的翻译，清代奏议按数字文本逐字引用。

\section{当时的世界}\label{sec:toku-world}

这一节交代：18 世纪初的日本与清代中国各有怎样的人口与赋税；两国的钱由谁发行，怎样计价。

\subsubsection{三千万人与“天下的厨房”}
按 Schaede 引用的估计，德川时代日本的人口“相当稳定，约为 3,000 万人”，其中 87\% 是农民，5\% 是武士，8\% 是住在城市里的商人与工匠\cite{schaede1989}。大阪是全国的商业中心，人口 1609 年约 20 万，到 18 世纪 50 年代翻了一番，达到 40 万，号称“天下的厨房”\cite{schaede1989}。土地按能产多少石米来计量，赋税以米征收，武士的俸禄也以米发放；1 石等于 10 斗，约 180 升\cite{schaede1989}。不过，Schaede 说，只有在德川时代的最初，米才被用作支付手段；“按这些标准，米从来不是严格意义上的货币”，可它是“基本的记账单位，尤其是在财政上”\cite{schaede1989}。各藩把收来的年贡米运到大阪，存进自己的仓库（藏屋敷），在那里卖成银子：这样的仓库，1673 年有 91 座，1730 年有 124 座\cite{schaede1989}。

\subsubsection{三种金属，两种计价}
幕府“几乎拥有全部贵金属矿藏，经营金、银、铜三种铸币厂”\cite{suzuki2021}。按 Suzuki 的说法，德川的货币制度“名义上是三本位的”：面额大的是金币，中等的是银币，小额的是铜钱；可是，“西日本主要的交换媒介是银与铜钱，东日本则是金与铜钱”\cite{suzuki2021}。两边的算法也不一样。在江户，货物“以金的记账单位标价”，按金币的面值计数，进位是四进制：1 两等于 4 分（bu），1 分等于 4 朱（shu）\cite{suzuki2021}。在西日本，货物“以银的记账单位标价，价格写成银币的重量，银币没有面值”：1 贯（kan）等于 1,000 匁（momme），一块丁银约重 161.25 克\cite{suzuki2021}。铜钱按枚计数，1 贯文等于 1,000 文\cite{suzuki2021}。Suzuki 说，东西之间的金银比价“与国际汇率相似”\cite{suzuki2021}：大阪与江户之间做生意，有点像两个国家之间做生意。

\subsubsection{清：称量的银与计数的钱}
清朝的两种钱，地位与日本不同。朝廷只铸铜钱。白银是法定的支付手段，但 Irigoin 说，它“从来不是政府以垄断权发行的货币”；“两”也没有统一的标准，各地的两“重量与成色各不相同”\cite{irigoin2013}。白银怎样在明代成为赋税的货币，见\chref{ch:silver}。银与钱之间有一个官定的比价。《清史稿》记载，顺治二年（1645）“定钱七枚准银一分”，即每两 700 文；“民间颇病钱贵，已更定十枚准一分”；顺治十年（1653），“钱千准银一两，定为画一通行之制”，钱粮“徵银七钱三”，即七成收银、三成收钱\cite{qingshigao}。Irigoin 也说，两种金属的比价“传统上在每单位白银 1,000 文铜钱上下”\cite{irigoin2013}。

市场的比价却很少正好是 1,000。18 世纪的麻烦是钱贵。康熙年间朝廷规定，“银一两易钱毋得不足一千，然钱直终不能平”，到康熙“季年银一两易钱八百八十至七百七十”；乾隆十三年（1748）以前，“京师钱价昂，银一两仅易八百文”，朝廷只好拿出工部库存的钱来平抑\cite{qingshigao}。乾隆在上谕里抱怨，“物之定直以银不以钱，而官民乃皆便钱不便银”\cite{qingshigao}。一百年后，吴嘉宾回顾说，“雍正以来。迄今几百年。钱价总不过一千一百文内外易银一两”\cite{xubiansheng}。到了 19 世纪，方向倒了过来，这是\secref{sec:toku-qing}的题目。

\section{石高与三货：以米计税}\label{sec:toku-koku}

这一节回答：以米计税与以银计税，谁承担米价与比价变动的风险？东边用金、西边用银，怎样抬高了江户的物价？

\subsubsection{谁承担米价的风险}
设一个纳税人的税额是 $T$ 两银子，他卖米得到铜钱，再用铜钱换银子交税。一两银子能换 $e$ 文铜钱，一石米能卖 $p$ 文钱，他为交税要卖的米就是
\begin{keyeq}
x=\frac{T\,e}{p}\quad\text{（石）} .\label{eq:toku-burden}
\end{keyeq}
税额 $T$ 不变，$e$ 翻一番而 $p$ 不变，要卖的米就翻一番；这正是曾国藩说的“暗加一倍之赋”。清代的田赋是这样计算的：包世臣（1775—1855）说，“各省正供。年额四千万两”\cite{xubiansheng}。以米计税的日本则不同：纳税人交的就是米，式~\eqref{eq:toku-burden} 里的 $x$ 是固定的，$e$ 与 $p$ 变动的风险落在收税的一方。一个藩在大阪卖 1 万石米，按岩桥整理的大阪米价，1721 年能收进 721 贯银子，1730 年只有 310 贯，少了 57\%\cite{iwahashi1981}。武士的俸禄以米计，可他们买的东西以金银计价；米相对于别的东西便宜了，武士就穷了。

\begin{intuition}[“米将军”为什么想要米贵]
德川吉宗（1716—1745 年在位）被称为“米将军”。Schaede 说，享保年间（1716—1735），幕府的米价政策来了一个“180 度的转弯”：它想要的“是降低一般物价，同时提高米价”\cite{schaede1989}。一个靠卖米过日子的政府，要的是米相对于其他货物变贵；一个按银两收税的政府，在银贵的时候不必着急，因为同样的税额能买到更多的东西，着急的是纳税人。可这只是账面上的好处。曾国藩说，“自道光十五年以前。江苏尚办全漕。自十六年至今。岁岁报歉。年年蠲缓”\cite{xubiansheng}：税额不变，收不上来的却越来越多。
\end{intuition}

\begin{compare}{以米计税与以银计税}{以米计税（德川日本）}{以银计税（清代中国）}
\cmprow{税额以什么计}{石米：土地按产米的石数计量}{银两：田赋“正供”年额 4,000 万两}
\cmprow{纳税人交什么}{米，由各藩运到大阪卖成银子}{不通商的地方交铜钱，由官府折成银两}
\cmprow{米价下跌时谁吃亏}{领主与拿米俸禄的武士}{纳税人：要卖更多的米}
\cmprow{比价上升时谁吃亏}{江户买西边货物的人（见下文）}{纳税人：银贵钱贱，同样的税要更多的钱}
\cmprow{政府希望什么}{米贵（1730 年的法令把“提高米价”写成第一条）}{钱价平稳（“银一两易钱毋得不足一千”）}
\cmpsum{以实物计税，价格风险落在收税的一方；以货币计税，价格与比价的风险落在纳税的一方。}
\end{compare}

\subsubsection{东边的金，西边的银}
江户卖的许多货物是从西日本运来的。Suzuki 说，江户是 17 世纪以来才逐渐发展起来的新城市，精致的货物大都从西日本的老城市运进来，大阪是西日本的商业中心\cite{suzuki2021}。同一件货，在大阪以银计价，在江户以金计价。设 $a$ 与 $b$ 是两种钱，一单位 $b$ 折合 $e$ 单位 $a$，同一件货以两种钱计的价格 $P_a$、$P_b$ 满足
\begin{keyeq}
P_a=e\,P_b,\qquad g_{P_a}\approx g_e+g_{P_b} .\label{eq:toku-twoprice}
\end{keyeq}
第二个式子取了增长率：以 $a$ 计的价格的涨幅，约等于比价的涨幅加上以 $b$ 计的价格的涨幅。在日本，$a$ 是银（匁），$b$ 是金（两）；在清代中国，$a$ 是铜钱（文），$b$ 是银（两）。把 $a$ 看作本地的钱、$b$ 看作外来的钱，这里的 $e$ 就与\appref{app:symbols}对汇率的约定（一单位外币折合多少本币）一致。

Suzuki 举了一个虚构的例子：如果大阪 1 两金换 40 匁银，江户 1 两换 80 匁，“西边的商人就把货物以西边两倍的价格运到江户。货物在江户用金币买下，然后把收到的金子在大阪换成银子，账就这样结清了”\cite{suzuki2021}。元禄改铸以前，东西两地的比价“极为稳定”，都是 60 匁换 1 两；改铸以后，比价剧烈波动，银子在大阪通常比在江户值钱，“这只在江户引起了通胀”\cite{suzuki2021}。江户的十组问屋（Tokumi Donya）因此苦恼，屡次请愿；1700 年，幕府规定全国一律以 60 匁银换 1 两金\cite{suzuki2021}。可这条法令“只有幕府自己与江户的商品市场遵守”；Suzuki 在脚注里说，东西两地货币市场上的浮动比价“相差不太大”，江户商品市场的 60 匁“不是真实货币的比价，而是根据西边货物的银价计算江户金价时用的比率”\cite{suzuki2021}。西边的商人“不遵守固定比价的法令”，按西边的市场比价定价\cite{suzuki2021}。

\begin{example}{东西两地的标价}\label{ex:toku-eastwest}
（1）用 Suzuki 的虚构例子：一件货在大阪值 40 匁银，大阪 1 两换 40 匁，江户 1 两换 80 匁。大阪商人要在江户收多少两金子，才能在大阪换回 40 匁？按江户的比价，这些金子值多少匁？（2）Suzuki 举了一个真实的例子：1718 年 11 月，商品市场上西边的比价是 1 两换 45 匁，东边是 60 匁\cite{suzuki2021}。一件在大阪值 45 匁的货，在江户要卖多少两？江户商品市场按 60 匁折算，标价是多少匁？（3）如果西边的比价从 45 匁变成 60 匁，江户的标价怎样变？
\solution （1）由式~\eqref{eq:toku-twoprice}，$P_b=P_a/e=40/40=1$ 两；按江户的比价，1 两值 80 匁，是大阪价格的两倍，这就是 Suzuki 说的“两倍的价格”。（2）$45/45=1$ 两；江户按 60 匁折算，标价是 60 匁，比大阪贵 $60/45-1\approx33\%$。一般地，江户标价与大阪价格之比是 $\bar e/e_W$，$\bar e=60$ 是江户的折算比率，$e_W$ 是西边的市场比价。（3）$e_W=60$ 时，$\bar e/e_W=1$，加价消失。西边的比价每变一次，江户的标价就得跟着改一次。比价是否可能长期相差这么多，可以用\chref{ch:silver}的式~\eqref{eq:silver-arbitrage} 检验：如果两地货币市场的比价相差 33\%，远超运送金银的成本，套利就会把它抹平；Suzuki 的脚注正是说，货币市场的比价相差不大，差别在江户商品市场那个固定的折算比率上。
\end{example}

\subsubsection{越后屋的换算表}
标价跟着比价变，代价不小。Suzuki 说，西边的浮动比价波动剧烈时，西边的商人“不得不大幅改动价签，这意味着巨大的菜单成本”\cite{suzuki2021}。和服店越后屋（今天的三越）的招牌卖法是“现金掛值无”（Genkin Kakene Nashi），Suzuki 解释为“用现金买和服，就按最低的价格卖”；这个办法靠的是写着最低价的价签，币制一乱，它就因为“巨大的菜单成本”陷入了危机\cite{suzuki2021}。1719 年，越后屋想出了解决办法：价签上写一个“加倍的价格”，另备一张以江户 1 两兑 60 匁的固定比率命名的“60 匁换算表”。店员要知道顾客选的和服的加倍价格、顾客打算怎样付款、西边的浮动比价，以及换算表上相应的比率；“把加倍的价格乘上相应的比率，店员就能立刻算出售价”\cite{suzuki2021}。这样，“即使西边的比价剧烈波动”，也不必改动价签\cite{suzuki2021}。今天的经济学把改价的成本叫作\textbf{菜单成本}（menu cost），它是新凯恩斯模型解释价格粘性的一个理由（\chref{ch:newkeynesian}）；越后屋的办法是把会变的部分从价签上拿掉，放进一张表里。

\section{元禄改铸与新井白石}\label{sec:toku-genroku}

这一节回答：幕府为什么改铸金银币？改铸以后物价为什么上涨，是因为钱多了，还是因为钱变差了？

\subsubsection{空了的金库}
Suzuki 说，德川初年贵金属产量丰富，币制运转良好；“没有人知道，金银资源几年之后就会几乎枯竭”。钱币成色好，人们用了一个世纪，可是荻生徂徕（1666—1728）说，幕府的金库“在一个世纪里就空了”\cite{suzuki2021}。改铸是重新利用有限资源、即市面上已有的钱币来充实金库的办法，“于是，幕府实行了 1695 年的元禄改铸”\cite{suzuki2021}。按室鸠巢（1658—1734）的说法，幕府的打算很简单：如果把国内的全部货币收回，“用收回的钱铸出两倍的数量，把其中一半投回流通”，幕府就能用同样多的钱把金库补满；Suzuki 说，这个计划“与钱的成色毫不相干”\cite{suzuki2021}。表~\ref{tab:toku-coins} 列出了 1601—1736 年的主要金银币。

\begin{table}[htbp]\centering\small
\caption{德川日本的小判与丁银，1601—1736 年}\label{tab:toku-coins}
\begin{tabularx}{\linewidth}{@{}l l Y Y Y Y@{}}
\toprule
金币（小判，面值 1 两） & 铸造年代 & 含金（\%） & 重量（克） & 纯金（克） & 铸造量（两） \\\midrule
庆长小判 & 1601—1695 & 84.291 & 17.85 & 15.05 & 14,727,055 \\
元禄小判 & 1695—1710 & 57.366 & 17.85 & 10.24 & 13,936,220 \\
乾字小判（宝永） & 1710—1714 & 84.291 & 9.38 & 7.91 & 11,515,500 \\
正德小判 & 1714 & 84.291 & 17.85 & 15.05 & 213,500 \\
享保小判 & 1716—1736 & 86.785 & 17.85 & 15.49 & 8,280,000 \\
元文小判 & 1736—1818 & 65.714 & 13.13 & 8.63 & 17,435,711 \\\midrule
银币（丁银，按重量计） & 铸造年代 & 含银（\%） & & & 铸造量（贯） \\\midrule
庆长银 & 1601—1695 & 80 & & & 1,200,000 \\
元禄银 & 1695—1706 & 64 & & & 405,850 \\
宝永二宝银 & 1706—1710 & 50 & & & 278,130 \\
永字银 & 1710 & 40 & & & 5,836 \\
三宝银 & 1710—1711 & 32 & & & 370,487 \\
四宝银 & 1711—1712 & 20 & & & 401,240 \\
享保银 & 1714—1736 & 80 & & & 331,420 \\
元文银 & 1736—1818 & 46 & & & 525,465.9 \\
\bottomrule
\end{tabularx}
\datasource{Suzuki（2021）表 2、表 3（据三井 1933 年与塚本 1923 年），按页面图像读出；纯金 = 重量 × 含金，由本书计算。Suzuki 注明表中庆长小判的铸造量偏低。}
\end{table}

\subsubsection{成色与物价}
改铸以后物价为什么上涨？当时城里常常议论这个问题，Suzuki 说，“人们把通胀归因于钱币成色的降低，而不是数量的增加”\cite{suzuki2021}。把这个看法写成式子：设一两钱币所含的纯金为 $c$ 克，一石米值 $q$ 克纯金。如果人们按钱币所含的金属给它估价，米以钱币计的价格就是
\begin{keyeq}
P=\frac{q}{c},\qquad \pi\approx g_q-g_c .\label{eq:toku-content}
\end{keyeq}
第二个式子取了增长率：物价的涨幅约等于米的金属价格的涨幅减去成色（每两的纯金）的增长率。改铸只动了 $c$，没有动 $q$；按这个式子，庆长小判换成元禄小判，每两的纯金从 15.05 克降到 10.24 克，以两计的物价就要乘上 $15.05/10.24\approx1.469$，涨 46.9\%。银币从含银 80\% 降到 64\%，以匁计的物价要乘上 $80/64=1.25$。

另一种解释是\chref{ch:silver}的数量方程，式~\eqref{eq:silver-qty}：改铸以后钱多了，$V$ 与 $Y$ 不变，物价就随 $M$ 上涨。Suzuki 说，今天的日本史学者“用货币数量论来解释元禄通胀”，可是“这样的理解是错误的”：新井白石收集了金银铸币厂官员的资料，按他的估计，因为许多旧币被窖藏起来，改铸期间“市面上流通的钱币数量没有变化”\cite{suzuki2021}。如果钱的数量没有变，数量论预测物价不变；而江户的米价涨了 46\%。

\begin{intuition}[钱为什么按成色打折]
式~\eqref{eq:toku-content} 成立，要有一条把钱的价值拴在金属上的绳子。德川日本有好几条。西边的银子本来就按重量计价，成色低了，同样重的银子就不值那么多。旧币可以留在家里不用，手里有好钱的人不会按面值把它花掉，这就是“劣币驱逐良币”；被留下来的旧币越多，市面上流通的就越是新币，物价也就越按新币的成色来定。商人也按成色给不同的钱作价（例~\ref{ex:toku-genbun}）。反过来，如果钱只是一种筹码，没有人能把它熔掉，也没有人在乎它含多少金子，物价就只取决于有多少筹码，这时才轮到数量论。
\end{intuition}

\subsubsection{改铸的收益}
幕府从改铸中得到多少？设收回旧币 $N$ 两，每两含纯金 $c_{\mathrm{old}}$，改铸成每两含 $c_{\mathrm{new}}$ 的新币，可以铸出 $N c_{\mathrm{old}}/c_{\mathrm{new}}$ 两；幕府按每两旧币付给持有人 $1+\phi$ 两新币，$\phi$ 是兑换的补贴率。不计铸造成本，幕府的收益是
\begin{keyeq}
S=N\left(\frac{c_{\mathrm{old}}}{c_{\mathrm{new}}}-1-\phi\right) .\label{eq:toku-seign}
\end{keyeq}
这是\chref{ch:medieval}的铸币税，式~\eqref{eq:medieval-seigniorage}，换成了日本的情形：一两旧币可以铸出 $ME=c_{\mathrm{old}}/c_{\mathrm{new}}$ 两新币，铸币厂付给送币人 $MP=1+\phi$ 两，铸币税率就是 $1-(1+\phi)/(c_{\mathrm{old}}/c_{\mathrm{new}})$。Suzuki 说，元禄改铸时，以旧换新给的补贴，到 1707 年为止“每 100 两庆长小判只有 1 两”，与成色的降低相比“微不足道”\cite{suzuki2021}。

\begin{example}{元禄改铸的账}\label{ex:toku-seign}
新井白石的报告说，8,824,350 两庆长金币与 287,617.155 贯庆长银币，“改铸成 13,236,534 两元禄金币与 357,535.380 贯元禄银币”；他又说，据铸币厂官员的说法，庆长金币共有 20,000,000 两，城里的庆长银币有 400,000 贯\cite{suzuki2021}。（1）每两旧金币铸出了多少两新金币？与表~\ref{tab:toku-coins} 的成色相比如何？银币呢？（2）按 1\% 的补贴，幕府从金币的改铸中得到多少两？铸币税率是多少？（3）金币与银币各有多大比例被改铸？如果没有改铸的旧币全被窖藏，市面上流通的金币是增加了还是减少了？如果全都照常流通呢？
\solution （1）$13{,}236{,}534/8{,}824{,}350\approx1.500$；按成色应为 $15.05/10.24\approx1.469$，报告的数字略高。银币：$357{,}535.380/287{,}617.155\approx1.243$，按成色应为 $80/64=1.25$，相当接近。（2）由式~\eqref{eq:toku-seign}，用实际铸出的数：$S=13{,}236{,}534-8{,}824{,}350\times1.01\approx432$ 万两；铸币税率 $1-1.01/1.500\approx32.7\%$。按成色计算则是 $1-1.01/1.469\approx31.3\%$。这正是室鸠巢说的“铸出两倍、投回一半”的简化版：每收回 100 两，幕府留下约 49 两，付出 101 两。（3）金币 $8{,}824{,}350/20{,}000{,}000\approx44\%$，银币 $287{,}617/400{,}000\approx72\%$；Suzuki 把它概括为“大体上可以说，一半的庆长币被改铸，其余被窖藏”\cite{suzuki2021}。如果没改铸的 1,117.6 万两全被窖藏，流通的金币从 2,000 万两降到 1,323.7 万两，少了 34\%；如果全都照常流通，则增加到 $1{,}323.7+1{,}117.6=2{,}441.2$ 万两，多了 22\%。Suzuki 的结论是流通的钱“没有变化或略有减少”\cite{suzuki2021}，介于两者之间。无论哪种情形，数量论都解释不了 46\% 的涨价。
\end{example}

\subsubsection{越改越差}
改铸一开了头，就停不下来。按 Suzuki 的表 3，银币的含银量 1706 年降到 50\%（宝永二宝银），1710 年降到 40\% 与 32\%，1711 年降到 20\%（四宝银）；金币在 1710 年改铸成乾字小判，成色恢复到庆长的 84.291\%，重量却只有 9.38 克，每两的纯金只剩 7.91 克（表~\ref{tab:toku-coins}）\cite{suzuki2021}。金与银贬值的幅度不同，金银比价就乱了。Suzuki 说，幕府从元禄改铸中学到，“金银币成色降低的比例不一致，会引起比价的失调”，后来的元文改铸才让两者保持了平衡\cite{suzuki2021}。改铸期间，幕府还想多向商人征税，可是“商人把税加进了商品的价格”\cite{suzuki2021}。

\subsubsection{新井白石的复原}
新井白石（Arai Hakuseki，1657—1725）是将军的顾问，曾上书反对改铸\cite{suzuki2021}。1714 年，幕府铸造了仿照庆长小判的正德小判，同年开铸含银回到 80\% 的享保银；1716 年又铸造了享保小判（表~\ref{tab:toku-coins}）\cite{suzuki2021}。新井白石相信，新铸的好银币能把窖藏的好银币引回流通。Suzuki 说，“他在银币上错了，在金币上却对了”：约 1716 年，过去二十年间从市面上完全消失的庆长小判又开始流通，按太宰春台（1680—1747）的说法，最终“占到了市面上小判的十分之九”\cite{suzuki2021}。

新旧钱币混在一起，比价就成了问题。1714 年，幕府采纳堺的商人谷长右卫门（Tani Chōemon）的提议，用法令规定了比价：庆长、享保小判 1 两等于乾字小判 2 两，元禄小判 1 两等于乾字小判 1 两\cite{suzuki2021}。银币的规定更能说明问题。此前，银币铸币厂用 1,500 匁四宝银收购 1,000 匁庆长银；1714 年以后，要付 2,000 匁；Suzuki 说，“按银的成色，应付 4,000 匁”，差额 2,000 匁由幕府承担\cite{suzuki2021}。1718 年的享保法令才“按成色规定了正确的比价”：庆长银、享保银 10,000 匁等于四宝银 40,000 匁\cite{suzuki2021}。元禄小判在 1718 年、乾字小判在 1720 年被禁止在市面上使用\cite{suzuki2021}。Suzuki 说，“好的坏的各种钱币都在市面上流通”；为了应付这种混乱，市面上还有人出售换算图表\cite{suzuki2021}。

钱变好了，以钱计的物价就跌了。按岩桥整理的大阪米价，1718 年每石 127.9 匁，1719 年只有 40.6 匁\cite{iwahashi1981}。Schaede 说，享保年间的米价“不仅名义上下跌，因为一连几年丰收，实际上也下跌了”\cite{schaede1989}。对一个以米计税的政府来说，这是坏消息。

\subsubsection{元文改铸}
1736 年，幕府又一次改铸。Suzuki 说，从享保小判到元文小判，纯金减少了 44.34\%，元禄改铸是 31.91\%；银币的纯银减少了 42.5\%，元禄是 20\%\cite{suzuki2021}。这一次，法令一面规定新的元文小判在市面上与旧的庆长小判“同样值钱”，幕府坚持“1 两就是 1 两，不论金币的成色”；一面又规定，拿 100 两旧小判到铸币厂，可以换到 165 两新小判，“10,000 匁旧银换 15,000 匁新银”\cite{suzuki2021}。Suzuki 称这是一种“既是名目论、又是金属论的矛盾政策”。它的一个目的是减轻武士的债务：改铸之后没几天，“许多大名就用新小判偿还了他们多不胜数的债务”；1736 年 5 月 28 日，江户的十组问屋向幕府申诉，按面值用新币还旧币的债，“债主会蒙受相当的损失”\cite{suzuki2021}。

\begin{example}{越后屋的大减价}\label{ex:toku-genbun}
越后屋事先得到了改铸的消息，从 1736 年 6 月 15 日起在江户与大阪大减价。Suzuki 记下了一件和服的价格：改铸以前，“加倍的价格是 120 匁（原价 60 匁 = 1 两），降到 48 匁，也就是 0.8 两，打八折”；用新小判付款，价格是 79.2 匁，即 1.32 两；新的 1 两小判含纯金 8.581726 克，而用旧小判付 0.8 两，相当于付出 11.9759232 克纯金\cite{suzuki2021}。（1）用旧币付款与用新币付款，价格之比是多少？与铸币厂的兑换率相比如何？（2）按表~\ref{tab:toku-coins}，1 两庆长小判与 1 两元文小判的纯金之比是多少？铸币厂的兑换率给持币人的补贴是多还是少？（3）谁从这次减价里得到了好处？
\solution （1）$1.32/0.8=1.65$，正是铸币厂 100 两旧币换 165 两新币的比率：越后屋按铸币厂的兑换率给旧币作价，没有按法令的“1 两就是 1 两”。（2）$15.05/8.63\approx1.744$；按式~\eqref{eq:toku-seign}，铸币税率 $1-1.65/1.744\approx5.4\%$，按享保小判算是 $1-1.65/1.795\approx8.1\%$，远低于元禄改铸的约 32\%。元文改铸给的是 Suzuki 所说的“适当的补贴”。（3）Suzuki 说，减价很受顾客欢迎，越后屋“每天挤满了远道而来的客人”，人们“愿意用旧小判买和服”\cite{suzuki2021}。按面值看，顾客用 0.8 两买到了原价 1 两的和服；按金子看，他们付出的 11.98 克，比用新币付 1.32 两的 $1.32\times8.58\approx11.33$ 克还多约 5.7\%，差额正是铸币厂扣下的那一部分。越后屋给旧币的作价与铸币厂一样，顾客用旧币付款，等于省去了自己去铸币厂兑换的手续；越后屋则收进了含金多的旧币。Suzuki 说，改铸的时候人们通常因为不安而不买东西，越后屋把这场危机变成了做生意的好机会\cite{suzuki2021}。法令说新旧等值，商店与顾客却都按金属给钱作价。
\end{example}

\begin{compare}{数量说与成色说}{数量说}{成色说}
\cmprow{主张}{物价随钱的数量变化，式~\eqref{eq:silver-qty}}{物价随每单位钱所含的金属变化，式~\eqref{eq:toku-content}}
\cmprow{当时与后来谁持此说}{今天的日本史学者；《管子》的轻重论}{江户城里“最普遍的意见”；Suzuki 据新井白石的估计}
\cmprow{对元禄改铸的预测}{流通量不变时物价不变，最多涨约 22\%}{江户米价涨 46.9\%}
\cmprow{成立的条件}{钱是筹码，熔不掉，也没人在乎它含多少金属}{钱可以熔化、窖藏，按成色兑换}
\cmprow{德川日本的例子}{幕府的法令：“1 两就是 1 两”}{西边按重量计价的银子；越后屋给旧币的作价}
\cmpsum{两种说法不是互相排斥的定律，而是两种极端情形；钱越接近按金属估价的商品货币，成色说越准，越接近纯粹的筹码，数量说越准。}
\end{compare}

Suzuki 引了 von Glahn 对《管子》的解读：钱币与货物的价值“成反比”，“可以看作一种货币数量论”；按 Suzuki 引的英译，《管子》的轻重论说，钱被看重，货物就都便宜，钱不值钱，货物就都贵\cite{suzuki2021}。Suzuki 认为，中国史书讲的几乎都是铜钱，铜钱越铸越薄、越铸越多，“成色与数量是成比例的”，所以中国的经验分不清两者；元禄改铸也是“成色更低、数量更多”，要靠新井白石对流通量的估计，才能把两种说法区分开\cite{suzuki2021}。\secref{sec:toku-test} 用米价检验它们。

\section{堂岛米市}\label{sec:toku-dojima}

这一节回答：大阪的米票怎样从仓库的提单变成期货合约？堂岛的交易制度怎样保证合约的履行？市场的价格有没有效率？

\subsubsection{米票}
各藩在大阪的仓库通过拍卖卖米，买主拿到的不是米，而是米票（米切手）：Schaede 说，仓库“发出一张对仓库里一定数量的米的所有权凭证，换取货币”\cite{schaede1989}。米票降低了在短时间内买卖大量米的成本。经纪人在米票到期以前就转手买卖，期限很快延长到 18 个月；仓库也不在乎，“开始发行没有米作担保的米票，而且不收保管费”\cite{schaede1989}。约 1700 年，米票被标准化为 10 石一张\cite{schaede1989}。

米票有几种。米到了仓库、拍卖之后付清了价款才发出的，叫“出切手”；没有存米就发出的，不写买主的名字与日期，叫“坊主切手”（monk bill），“所以坊主切手类似于无记名债券”\cite{schaede1989}。还有一种不经拍卖、凭预付款发出的“空米切手”，Schaede 说，“发行人（仓库）与购买人都不把它们看作实物米的所有权凭证”，它们更像商人借给仓库的一笔钱的凭据\cite{schaede1989}。他把区别概括为：“买主如果在买进一张没有担保的米票的当天付清全部款项，他实际上买的是一张债券；反之，如果米商买的是有担保的米票，他就订立了一份远期合约。”\cite{schaede1989}没有担保的米票很多。按一本米价笔记，1749 年流通中的米票代表了 11 万多俵米，而当时仓库里的存米只有约 3 万俵，“几乎是可供实物交割的米的四倍”\cite{schaede1989}。Schaede 说，常年缺钱的大名在大阪有两种融资办法，一是直接向大阪商人借款，二是“发行‘没有担保’的米票，即交割米的空头承诺”\cite{schaede1989}。

\subsubsection{从禁止到承认}
幕府起初把米票当作物价的祸根。Schaede 说，幕府“不论对错，都认为没有担保的米票的膨胀是通胀的主要原因”：1652 年禁止没有担保的米票，1660 年干脆禁止米票交易，并把米票的期限限制在 30 天以内，1663 年又缩短到 10 天\cite{schaede1989}。禁令没有效果：“商人们现在在 10 天之内付清米票的全部款项，米仍旧留在仓库里，米票照样活跃地交易”\cite{schaede1989}。

\begin{case}{大阪 · 1705—1730 · 从查封淀屋到承认堂岛}
大阪最早的米市开在淀屋（Yodoya）的门前，Schaede 称淀屋是“当时首屈一指的商家”。淀屋在通往北浜的大路上，每天聚集在那里的米商“妨碍了交通”，1688 年，官府让米商们集中到堂岛，一个三条河交汇的河口上的小岛；1697 年，淀屋自己也搬到了堂岛\cite{schaede1989}。幕府这时把物价上涨归咎于“帐合米”交易，即在淀屋门前进行、被说成“不过是赌博”的期货交易。1705 年，“淀屋被查封，它对大名的债权被宣布无效，它的巨额财富被没收”；官方的理由是淀屋违反了禁奢令，Schaede 说，“真正的原因是希望取缔‘虚假的’赌博能阻止米价上涨”\cite{schaede1989}。交易照样在淀屋被查封的门前暗中进行。

二十多年以后，问题反了过来：米价太低。Schaede 说，幕府认为堂岛米商的交易是“虚假的”“抬高价格的”，于是在 1728 年正式准许延期交易，1730 年，堂岛米市成为日本“唯一得到官方承认、有组织的期货交易所”\cite{schaede1989}。1730 年的法令写明：（1）正式准许这个市场的目的，“是提高米价”；（3）清算业务只限于此前已在做这一行的 50 家清算所；（5）只有堂岛可以做帐合米交易\cite{schaede1989}。幕府本来只打算临时准许，“一旦米价回升，就禁止帐合米交易”；但是，“市场一旦正式建立起来，幕府就关不掉它了”\cite{schaede1989}。同一个市场，米价高时被当作抬价的赌场查封，米价低时又因为能抬价而被承认：幕府两次都相信期货交易能左右现货的价格。
\end{case}

\subsubsection{交易的结构}
承认以后，幕府出售交易许可（株）来控制米商的人数：1731 年 500 张，1732 年 4 月又 500 张，11 月 300 张\cite{schaede1989}，共 1,300 名有执照的米商。清算所（米方两替）1731 年领到 50 张特许，1746 年又增加 10 张；它们“不许为自己交易，只能执行委托”，“收到保证金以后，就承担合约履行的责任”\cite{schaede1989}。清算所不能为逐日清算收费，Schaede 说，“大概因为逐日清算是一件极费人力的生意”，它们在这上面赔钱；最初的六十家，只有四家活到了 19 世纪\cite{schaede1989}。图~\ref{fig:toku_dojima} 画出了 Schaede 描述的流程。拍卖时，问屋（tonya）先付 30\% 作为保证金，拿到一张写着总价的“银切手”，“在此后 30 天内付清 70\% 的余款，拿到米票”（第一阶段，远期）\cite{schaede1989}；米票可以在堂岛的市场上卖出（第二阶段，远期或期货）；最后，想要实物米的批发商买下米票，凭票提米（第三阶段）。

\fig{toku_dojima}{堂岛米市的交易流程。看三个阶段怎样把一张提单变成期货合约：第一阶段，问屋在仓库的拍卖上付 30\% 的保证金，拿到银切手，30 天内付清 70\% 的余款，拿到米票，这是远期交易；第二阶段，米票与以 100 石为一张的标准合约在堂岛交易，清算所收取至多三分之一的保证金，每十天结算一次盈亏，并保证合约的履行，这是期货交易；第三阶段，需要实物米的批发商买下米票提米。示意图，按 Schaede（1989）图 1 与第 4、5 节重画。}{fig:toku_dojima}

帐合米，即“账面上的米”的交易，就是期货。按 Schaede 的描述，一年分成三个交易期，两期之间休市约 10 天，“以免把未平的头寸展期到下一期”\cite{schaede1989}。一张合约“写的是 100 石米”，交易的是各期指定的一种标准米（建物米）；报价按合约张数，“‘3 张，23’就是 300 石米，每石 23 匁”\cite{schaede1989}。建立期货头寸，要在清算所存入“至多相当于交易总额三分之一”的保证金；与今天不同，保证金“随合约一起转手，由下一个买主接过去”\cite{schaede1989}。收盘价的定法别具一格：交易所的房梁上挂一个装着火绳的木盒，点燃火绳，火灭时的价格就是当天的正式收盘价，叫“火绳值段”；可是交易的人并不理会，要由“水方”泼水驱散人群，泼水不够，就“一桶一桶地把水浇在人群身上，这通常就停止了当天的交易”；这时的价格叫“桶值段”，“是记入账簿、用于逐日结算或交割的实际价格”\cite{schaede1989}。如果火绳烧完时没有成交价，当天的交易全部作废，Schaede 说，这样一来，“这种制度排除了囤积垄断”\cite{schaede1989}。

清算按十天一次进行。每晚米商把当天的交易抄进清算所的账簿；每十天，各清算所的书记在中央清算处会面，核对客户的头寸，“亏损的人要补交差额（维持保证金），盈利的人从清算中心领到款项”，Schaede 称之为“十天一次的逐日结算机制”\cite{schaede1989}。1737 年起，每期最后三天可以用现银或现米结清全部头寸；但实际上，“除了德川时代的最末期，用实物交割的合约很少”\cite{schaede1989}。

\subsubsection{套期保值}
期货的用处之一是转移价格风险。Schaede 说，持有米票的人可以“同时建立一个期货头寸，对冲米票资产可能的价格下跌”\cite{schaede1989}。设一个米商持有 $Q$ 石米的米票，同时按价格 $F_0$ 卖出 $Q$ 石的期货。到期时现货价格为 $S_T$，期货价格为 $F_T$，他的米票值 $QS_T$，期货头寸盈亏 $Q(F_0-F_T)$，合计
\begin{keyeq}
V_T=QS_T+Q(F_0-F_T)=QF_0+Q(S_T-F_T),\qquad \frac{\text{保证金的盈亏}}{\text{保证金}}=\frac{\Delta F}{m\,F_0} .\label{eq:toku-hedge}
\end{keyeq}
第一个式子说明，只要到期时期货价格与现货价格一致（$F_T=S_T$），他就锁定了 $QF_0$，米价涨跌与他无关；剩下的风险是两者之差 $S_T-F_T$，叫基差。第二个式子说明保证金的杠杆：保证金是合约价值的 $m$ 倍，期货价格变动 $\Delta F/F_0$，保证金就变动 $\Delta F/(mF_0)$。Schaede 说，保证金“似乎在交易额的 30\% 上下波动”；1770 年降到每石“一到两匁”，他算作每张 100 石的合约 200 到 300 匁，“最低的保证金只有 5\%”\cite{schaede1989}。

\begin{example}{一张合约的对冲}\label{ex:toku-hedge}
教学假设：期货价格 $F_0=60$ 匁/石，与 1770 年大阪的米价（66.3 匁\cite{iwahashi1981}）相近。一个米商持有 100 石米的米票，卖出一张期货合约。（1）保证金按三分之一计，他要存多少？到期时米价跌了 10\%，期货与现货都是 54 匁，他的米票与期货各盈亏多少？合计多少？（2）米价在头十天就跌了 10\%，按十天一次的结算，他的期货头寸收到多少？对买进合约的一方来说，亏损占保证金的比例是多少？保证金如果是 1770 年的 5\% 呢？（3）如果到期时现货是 54 匁，期货是 55 匁呢？
\solution （1）保证金 $100\times60/3=2{,}000$ 匁。米票亏 $100\times(54-60)=-600$ 匁，期货赚 $100\times(60-54)=600$ 匁，合计 $V_T=100\times60=6{,}000$ 匁，正是锁定的价值。（2）十天以后结算，他从清算所领到 600 匁，买方要补交 600 匁。由式~\eqref{eq:toku-hedge}，亏损占保证金 $0.10/(1/3)=30\%$；保证金为 5\%（每张 300 匁）时，是 $0.10/0.05=200\%$，买方必须补交两倍于原保证金的钱，否则清算所就要承担违约。保证金越低，买方越可能付不出，清算所承担的风险就越大。（3）$V_T=6{,}000+100\times(54-55)=5{,}900$ 匁，基差让他少了 100 匁。
\end{example}

\subsubsection{市场有没有效率}
如果期货价格是到期现货价格的无偏预测，就应当有 $S_T=a+bF_0+u_T$，且 $a=0$、$b=1$。Takatsuki 用这个回归检验堂岛的市场，其中 $S_T$ 是交易期到期日的现货价格，$F_0$ 是交易期开始时的期货价格\cite{takatsuki2007}。他引 Ito（1993）的结论说，1763—1780 年，堂岛“没有满足无偏条件”；他自己的检验是，1798—1835 年，在 1\% 的显著性水平上不能拒绝无偏假设，1840—1864 年则被拒绝，那时堂岛“因为一些交易者的违法交易而频频停市”\cite{takatsuki2007}。按弱式有效的检验，现货市场约 80\% 的时期、期货市场约 70\% 的时期不能拒绝原假设\cite{takatsuki2007}。Schaede 还引了宫本又郎的研究：1751 年到 19 世纪 30 年代，现货与期货价格“高度相关”\cite{schaede1989}。

市场最终坏在它自己的收盘规则上。天保年间（1830—1843），米的现货与远期价格随物价上涨，“期货价格却没有涨，因为这套制度使正常的定价变得不可能”；到 1866 年 4 月，官府不得不关闭远期与帐合米的交易，1869 年明治维新以后，堂岛米市被关闭\cite{schaede1989}。1871 年，堂岛商品交易所开张\cite{schaede1989}。

\begin{compare}{堂岛米票与现代期货}{堂岛的帐合米（1730 年以后）}{现代期货市场}
\cmprow{合约}{标准化，每张 100 石，各期指定一种标准米}{标准化的数量与等级}
\cmprow{参与者}{1,300 名有执照的米商}{交易所会员与经纪商}
\cmprow{履约保证}{清算所收保证金后承担履约责任；中央清算处在违约时承担合约义务}{中央对手方清算}
\cmprow{保证金}{约 30\%，低时 5\%；随合约转手，不是客户自己的账户}{客户在经纪商处的个人账户，逐日补足}
\cmprow{结算}{十天一次（早期每天了结全部头寸）；到期以现银或现米结清，实物交割很少}{逐日结算；多数头寸在到期前平仓}
\cmprow{收盘价}{火绳与水桶；无成交则当日交易作废}{交易所按规则确定的结算价}
\cmpsum{堂岛已经有了标准合约、清算所担保、保证金与定期结算这些期货市场的要件；Schaede 指出的三点不同是逐日了结全部头寸、保证金随合约转手、收盘规则让交易者能左右收盘价是否存在。}
\end{compare}

\begin{pitfall}[期货市场只是赌场]
幕府在 1705 年这样看，在 1730 年又认为期货能抬高米价；Schaede 在脚注里说，美国中西部的农民也曾多年试图以违反反赌博法为由关闭芝加哥期货交易所\cite{schaede1989}。期货交易确实是零和的：一方的盈利就是另一方的亏损。但它的用处不在于创造财富，而在于把价格风险转给愿意承担的人（例~\ref{ex:toku-hedge}），并且每天公布一个人人都能看到的价格。Takatsuki 的检验说明，在交易制度运转正常的时期，不能拒绝“堂岛的期货价格是到期现货价格的无偏预测”这个假设；坏的时期，坏在定价规则与违法交易，而不在期货本身。
\end{pitfall}

\section{清代的银与钱}\label{sec:toku-qing}

这一节回答：19 世纪上半叶，一两银子能换的铜钱为什么从约 1,000 文涨到 2,000 文以上？这对纳税人与挣工资的人意味着什么？

\subsubsection{从钱贵到银贵}
18 世纪的钱贵，在 19 世纪变成了银贵。图~\ref{fig:qing_cash}(a) 的实线是 Gamble 整理的北京郊区的银钱比价：1807 年每两银子换 979 文钱，1824 年 1,208 文，1838 年 1,553 文，1846 年 2,010 文，1853 年 2,205 文\cite{gamble1943}。按\chref{ch:songming}用过的年均复合变化率，1807—1853 年，银价平均每年上涨约 1.8\%。当时人的说法与此相符。道光十八年（1838），鸿胪寺卿黄爵滋（1793—1853）上奏，据魏源（1794—1857）的转述，“向时纹银每两兑钱千今则每两兑至千有六百”\cite{xubiange}。包世臣在道光十九年（1839）记下，五月江西省城的银价“长至制钱一千。兑库纹六钱一分。是银每两为钱一千六百三十余文”\cite{xubiansheng}，$1000/0.61\approx1{,}639$。道光二十五年（1845），刘良驹上疏说，“京中纹银每两易制钱几及二千文。外省则每两易制钱二千二三百文不等”（数字文本的断句有误，这里按文意改正）\cite{xubiansheng}。吴嘉宾说，“现在半年之间。银价自一千五百文骤长至二千文有零”\cite{xubiansheng}。《清史稿》的总结是，道光十七年（1837）以后，“海内银卒耗竭，每两易钱常至二千”\cite{qingshigao}。

\fig{qing_cash}{清代的银钱比价与工资，1807—1853 年。(a) 看实线怎样从 1,000 文附近升到 2,000 文以上：实线是北京郊区一两银子折合的铜钱（Gamble 整理），褐色方块是正文引用的四个当时人的说法：黄爵滋（1838）、包世臣（1839）与刘良驹（1845）说的京中、外省两个数，外省的“二千二三百文”取中点。(b) 看两条线怎样分开：实线是北京郊区非熟练工人的日工资（铜钱），褐色虚线是同一工资折成的银两，都以 1807—1813 年的平均为 100；1845—1850 年，铜钱工资只高 4.8\%，银两工资低 45\%。数据：Gamble（1943），据 GPIH 的工作簿（工资已按 Allen 等人的处理加倍以计入伙食）；当时人的说法见《清史稿》、盛康与葛士浚的两种《皇朝经世文续编》。}{fig:qing_cash}

\subsubsection{两种钱，一个价格}
用式~\eqref{eq:toku-twoprice}，$a$ 是铜钱，$b$ 是银两：一件货以钱计的价格 $P_c$ 等于银钱比价 $e$ 乘以它以银计的价格 $P_s$，$g_{P_c}\approx g_e+g_{P_s}$。$e$ 涨了一倍，如果以钱计的价格不变，以银计的价格就要跌一半。曾国藩说的正是这种情形：石米“买钱三千。自古迄今。不甚悬远”，以银计的米价从三两跌到一两五钱\cite{xubiansheng}。包世臣说，这两三年收成好，收获时“谷一石。仅值钱五百上下”，青黄不接时“谷价仍不过七百数十文”\cite{xubiansheng}。以钱计的物价没有跟着银价涨，以银计的物价就跌了。Irigoin 概括传统的说法时也说，白银外流一方面引起了铜钱的通胀，另一方面使“以银计的价格暴跌，加重了某些以银征收的关键税种，尤其是田赋的负担”\cite{irigoin2013}。

工资也是这样。图~\ref{fig:qing_cash}(b) 是 Gamble 整理的北京郊区非熟练工人的日工资：1807—1813 年平均每天 81.25 文，1845—1850 年平均 85.2 文，只多了 4.8\%；同一时期一两银子从平均 1,036 文涨到 1,990 文，涨了 92\%；折成银两，日工资从 0.0784 两降到 0.0428 两，少了 45\%\cite{gamble1943}。按式~\eqref{eq:toku-twoprice}，$g_{w_s}\approx g_{w_c}-g_e$：铜钱工资几乎不动，银两工资就随比价下跌。对一个收铜钱、买以铜钱计价的东西的工人来说，这没有多大关系；有关系的是要用银子付的东西，比如赋税。

\begin{example}{银贵钱贱下的税负}\label{ex:toku-burden}
（1）曾国藩说，东南每石米卖钱三千，从前一两银换一千文，现在换两千文；从前卖米三斗就足够交一亩的税，现在卖六斗还不够\cite{xubiansheng}。用式~\eqref{eq:toku-burden}，三斗米对应的税额是多少两？比价翻番以后，同样的税要卖多少米？（2）包世臣说，江西新喻县“每钱粮一两。柜收花户钱一千八百八十五文”，谷价收获时每石五百文上下，青黄不接时七百数十文，“是小民完银一两。非粜谷二三石不可”\cite{xubiansheng}。按这两个谷价，完银一两各要卖多少谷？（3）吴嘉宾说，“卖十石谷。袛易三两银。谷每石六百。银每两值钱二千”，“佣一年工。袛易五两银。佣值岁不过十千”\cite{xubiansheng}。核对他的算术；如果比价回到一千一百文，同样的谷与工钱值多少银子？
\solution （1）$x=0.3$ 石，$p=3{,}000$，$e=1{,}000$，$T=xp/e=0.9$ 两；$e=2{,}000$ 时，$x=0.9\times2000/3000=0.6$ 石，正是曾国藩的“六斗”。（2）$1885/500\approx3.8$ 石；七百数十文按 750 文算，$1885/750\approx2.5$ 石；包世臣的“二三石”是按青黄不接时的谷价算的，收获时还要更多。注意，按“定例制钱一千。准库纹一两”，完银一两本应只要一千文\cite{xubiansheng}。（3）$10\times600/2000=3$ 两，$10{,}000/2000=5$ 两，算术无误。按 1,100 文，十石谷值 $6000/1100\approx5.5$ 两，一年的工钱值 $10000/1100\approx9.1$ 两：同样的谷与劳动，在银子上的价值几乎少了一半。
\end{example}

\section{银贵钱贱的争论}\label{sec:toku-debate}

这一节回答：银为什么贵了？是鸦片贸易让白银流出了中国，还是有别的原因？这个问题当时就有争论，今天的研究仍在争论。

\subsubsection{白银外流说}
最流行的解释是白银外流。《清史稿》说，道光年间“番船冒禁，岁漏出以千万计，御史黄中模、章沅咸以为言”；外国的银元流入，“至洋银日多，纹银日少而贵”\cite{qingshigao}。道光十八年，黄爵滋上奏说：“查道光初年，岁漏银数百万”，“十四年以前，岁漏二千馀万”，“近年岁漏三千馀万。此外各海口合之亦数千万”\cite{qingshigao}。魏源转述他的奏疏，把原因归于“粤东洋船鸦片烟盛行致纹银透漏出洋”\cite{xubiange}。林则徐（1785—1850）支持黄爵滋，并算了一笔账。

\begin{example}{林则徐的算术}\label{ex:toku-lin}
林则徐说：“吸鸦片者每日除衣食外至少亦须另费银一钱是每人每年即另费银三十六两”；“以户部历年所奏各直省民数计之总不止于四万万人若一百分之中仅有一分之人吸食鸦片则一年之漏即不止于万万两”；黄爵滋所说的“岁漏银数千万两尚系举其极少之数而言耳”。他还说，“若以食贫之人当中熟之岁大约一人有银四五分即可过一日”，“若一日有银一钱则诸凡宽裕矣”\cite{xubiange}。（1）一钱一天，一年按多少天算？（2）四亿人里百分之一吸食，一年共花多少银子？与包世臣说的“正供”四千万两相比如何？（3）一天一钱的鸦片，相当于一个穷人几天的生活费？（4）林则徐算出的是“漏银”吗？
\solution （1）一两等于十钱，$36\times10=360$ 天。（2）$4\times10^8\times1\%\times36=1.44$ 亿两，确实“不止于万万两”，是正供的 $1.44/0.4=3.6$ 倍。（3）一钱等于十分，$10/5=2$ 到 $10/4=2.5$ 天。（4）不是。这是吸食者在鸦片上的花费，不是净流出的白银：鸦片的价钱里有国内的转运与零售，进口的鸦片也可以用茶、丝等出口货物抵付，只有贸易差额的那一部分才需要用白银结清（\chref{ch:mercantilism}孟的贸易差额）。林则徐的算术说明了鸦片的花费之大，却不能证明每年有一亿多两白银流出。
\end{example}

\subsubsection{窖藏、用银与“钱贱非钱多”}
另一些人不相信流出能解释一切。刘良驹说：“今之银贵由于钱贱。实由于用银太重用钱太轻”；“纹银之少。亦未必非豪猾富商。利轻赍而储厚积”；他的办法是多用钱：“纹银日少。势不能使之更多。制钱日轻。法能使之渐重”；他建议由户部定一个银钱比价，像近年各省捐输案内“以制钱一千五百文抵银一两之例”那样，规定哪些收支可以收钱、放钱\cite{xubiansheng}。吴嘉宾说得更形象：说银贵是因为纹银出洋，“此犹饮水者忧天旱水涸。不知特釜中涸耳”；“天下之人。小者积以箱箧。大者积以瓮窖。人情之所私。安得不贵”；“今之钱贱。本非钱多。以上下皆便用银。富者又多藏银。始不敷用。北人行使空票。南人多用洋钱”\cite{xubiansheng}。在他看来，河里有水，是锅里干了：问题不在全国的白银少了，而在大家都想持有白银，不肯拿出来用。

还有人认为银贵根本无害。周腾虎说：“向者银足之时。一两仅值钱六七百。今者一两。率准钱千五六百。夫银多则折价以平钱。银少则昂价以适用。其用一也。银质虽去。银用仍留。何损于中国哉”；“天下之贫富在衣食而不在货币”\cite{xubiansheng}。这是货币中性的论证：钱的单位换了，实际的东西没有变。它错在忽略了以银两固定下来的债务与税额：田赋按银两计，工资与物价按铜钱计，比价一变，财富就在纳税人与收税人、债务人与债权人之间转移，这正是例~\ref{ex:toku-burden} 的算术。

\subsubsection{今天的研究}
今天的研究同样分成几派。Atwell 介绍林满红（Man-houng Lin）的《颠倒的中国》（\emph{China Upside Down}，2006）时说，她“不认同那种老看法”，即 19 世纪上半叶中国的白银外流“只是鸦片进口增加的结果”\cite{atwell2008}。按 Atwell 与 Irigoin 的转述，林满红强调供给的一面：中国货币制度的危机，“至少部分地是由拉丁美洲争取政治独立的长期斗争引起的”，这场斗争一再“打乱了［秘鲁、玻利维亚与墨西哥的］银矿与全球的白银供给”，1808 年以后流入中国的白银“急剧减少”；世界对生丝、丝织品与茶等中国产品的需求也急剧下降\cite{atwell2008,irigoin2013}。Atwell 还说，银钱比价的扩大，有一部分也是因为贵金属升值以后人们窖藏得更多；1850 年代中期以后，白银又“大量流回中国”\cite{atwell2008}。Irigoin 引 von Glahn 的看法，国内白银市场的失衡降低了对白银进口的需求\cite{irigoin2013}。

Irigoin 自己的解释着眼于钱币的质量。她认为，中国对白银的需求，到 19 世纪 20 年代为止，“主要是对一种可靠的货币标准的需求，而不是对银块的需求”：这个标准就是西班牙的卡洛斯银元（Carolus dollar）\cite{irigoin2013}。西属美洲独立以后，“中国人很快发现，19 世纪 10 年代开始出现的新银元靠不住。足值的旧币被窖藏起来，变得更加稀少”；卡洛斯银元供给的骤减“必然造成通缩冲击，而好钱被窖藏、流通速度下降，又加重了这一冲击”，这“也许可以部分解释道光年间的铜钱通胀”\cite{irigoin2013}。1856 年的报道说，在上海，卡洛斯银元比同样重的墨西哥银元或其他银元有 50\% 的升水\cite{irigoin2013}。这是元禄改铸的翻版：人们按质量给钱估价，好钱被窖藏，劣币留在市面上。她也评论了林满红的论证。按 Irigoin 的转述，鸦片进口在道光年间增加到原来的 20 到 25 倍；林满红把鸦片与白银的关系倒了过来，把白银的稀缺归于“传统来源的枯竭”，并用鸦片进口与白银外流“在时间上的相关”说明鸦片进口是对白银短缺的反应。Irigoin 说，“相关不等于因果”，“有意思的问题是，因果关系”究竟是从墨西哥的银矿到中国，还是反过来\cite{irigoin2013}。

时间顺序可以帮助判断。Irigoin 说，中国与英国的交易在 19 世纪 10 年代变为逆差，白银外流从 20 年代中期才开始\cite{irigoin2013}。按 Gamble 的数据，1807 年到 1824 年，银钱比价已经从 979 文涨到 1,208 文，涨了 23\%\cite{gamble1943}。如果这些时间大致不错，银价在外流开始以前就在上涨，单是鸦片与外流解释不了全部的银贵。

\begin{pitfall}[银贵钱贱只是白银流出的结果]
这是 1838 年黄爵滋与林则徐的看法，也是教科书里常见的说法。它把三件事混在了一起。第一，流出的白银有多少，当时人的数字差别很大，从“数百万”到“数千万”，林则徐的“万万两”是鸦片的花费，不是流出的数量（例~\ref{ex:toku-lin}）。第二，比价是两种钱的相对价格，铜钱一方的变化同样重要；刘良驹、吴嘉宾都认为问题在窖藏与“用银太重”，Irigoin 认为在好银元的短缺。第三，银价在外流开始以前就在上涨（\secref{sec:toku-test}）。鸦片贸易很可能加重了银贵，但“银贵钱贱”是供给、需求与钱币质量共同作用的结果。鸦片贸易与白银流动在 19 世纪后半叶的后果，见\chref{ch:periphery}。
\end{pitfall}

\begin{case}{北京 · 1853—1861 · 咸丰大钱与宝钞}
《清史稿》说，“先是道光中叶，银外泄而贵，朝野皆欲行大钱以救之”\cite{qingshigao}。咸丰初年，太平天国起事，财政告急：“是时银亏钱匮重，而军需河饷糜帑二千数百万，筹国计者，率以行官票请”\cite{qingshigao}。清朝发行纸币，不是第一次：《清史稿》说，“钞尝行于顺治八年，岁造十二万八千有奇。十年而罢”（“十年而罢”可以读作行了十年，也可以读作顺治十年停止）\cite{qingshigao}。宋元明纸币的兴衰，见\chref{ch:songming}。

这一次，朝廷同时发行了三种东西：以银两计的官票（“额题‘户部官票’”）、以制钱计的宝钞（“钞额题‘大清宝钞’”），以及面值“当千至当十，凡五等”的大钱；按规定，“票银一两抵制钱二千，钞二千抵银一两”\cite{qingshigao}。比价写在纸上，人们却不按纸上的比价用钱。当十大钱流通得最久，“然一钱当制钱二，出国门即不通行”\cite{qingshigao}：面值十文的钱只当两文用，出了北京城就没人收。宝钞也一样，“其后京师以官号七折钱发钞，直益低落，至减发亦穷应付，钞遂不能行矣”\cite{qingshigao}。这与元禄改铸有相同的地方：钱的面值抬高了，含的铜没有跟着增加，人们就按含铜量给它估价。不同的是，幕府垄断了金银矿与铸币厂，新币在全国通行；清朝的大钱出了京城就不通行。
\end{case}

\section{东亚与欧洲}\label{sec:toku-europe}

这一节回答：与\chref{ch:finrevolution}和\chref{ch:mercantilism}讲的同期欧洲相比，东亚的货币与金融制度有什么相同，有什么不同？

表~\ref{tab:toku-europe} 把三个地方放在一起。只列本章与前两章的史料能说明的方面。

\begin{table}[htbp]\centering\small
\caption{德川日本、清代中国与西欧，17 世纪末到 18 世纪}\label{tab:toku-europe}
\begin{tabularx}{\linewidth}{@{}l Y Y Y@{}}
\toprule
方面 & 德川日本 & 清代中国 & 西欧（荷兰、英格兰、法国） \\\midrule
谁铸钱 & 幕府垄断金银铜矿与铸币厂 & 朝廷只铸铜钱，白银由民间称量 & 君主的铸币厂；荷兰另有银行货币 \\
怎样计价 & 东边金币按面值，西边银子按重量 & 银两按重量与成色，铜钱按枚 & 记账单位（里弗尔、英镑）与铸币的法定作价 \\
货币贬值 & 1695—1711 年多次改铸，1714 年复原，1736 年再改铸 & 咸丰大钱（1853）；私铸小钱 & 1720 年法国金属里弗尔贬值三分之一 \\
期货与证券 & 堂岛米票与帐合米（1730） & —— & 公司股份与国债在阿姆斯特丹、伦敦、巴黎交易 \\
政府借债 & 大名向大阪商人借款，发行无担保米票；1773 年幕府为米票担保 & 1853 年发行官票、宝钞 & 永续年金、可转让国债与英格兰银行 \\
金融中介 & 两替商、挂屋 & “北人行使空票”的私票 & 阿姆斯特丹银行、英格兰银行 \\
\bottomrule
\end{tabularx}
\datasource{本章引用的 Suzuki（2021）、Schaede（1989）、《清史稿》与《皇朝经世文续编》；欧洲一栏据\chref{ch:finrevolution}与\chref{ch:mercantilism}。“——”表示本章的史料没有可比的制度。}
\end{table}

\subsubsection{相同之处}
三个地方的商人都在应付多种钱币与浮动的比价。江户的市面上有人卖换算图表，越后屋用 60 匁换算表，清代的商人要应付各地重量与成色各不相同的两（\secref{sec:toku-world}），阿姆斯特丹银行用银行货币统一记账（\chref{ch:finrevolution}）。三个地方的政府都用改动钱币的办法筹钱：元禄改铸、劳在 1720 年的贬值（\chref{ch:mercantilism}）、咸丰大钱。Suzuki 说，幕府铸出的钱“流进了商人控制的市场，再也没有回来”\cite{suzuki2021}。

\subsubsection{不同之处}
最大的不同在政府的债务。英格兰在 1690 年代以后靠议会担保的长期国债借钱（\chref{ch:finrevolution}）；德川日本没有这样的公债，缺钱的大名向大阪的商人借款，或者发行没有米作担保的米票。Schaede 说，两替商成了大名的财务代理人，凭将来运来的米向他们放款；仓库的财务代理人挂屋，生意“与现代的银行非常相似”\cite{schaede1989}。1773 年，幕府开始为全部米票担保，“实际上把各藩为赤字融资发的米票变成了幕府担保的债券”\cite{schaede1989}；幕府还规定，有人告米票欺诈或违约，“如果诉讼有理，政府就支付赔偿”\cite{schaede1989}。Schaede 说，有了这些担保，米票实际上相当于地方债券或幕府担保的债券，只是以米计值；它却没有发展成英格兰那样的国债市场。本章的清代史料里没有可比的长期公债，1853 年的军费靠的是官票、宝钞与大钱。另一个不同在期货：堂岛米市比 Schaede 所举的芝加哥期货交易所（1848）早了一个多世纪，后者之后还有法兰克福（1867）与伦敦（1877）\cite{schaede1989}。

\section{数据检验与反例}\label{sec:toku-test}

这一节回答：成色说能不能解释德川日本的米价？清代的银钱比价与工资，是否符合式~\eqref{eq:toku-twoprice} 的预测？

\subsubsection{江户的金价米}
用式~\eqref{eq:toku-content} 做一个简单的预测。以 1685—1694 年江户米价的平均 0.786 两/石为基准，那时流通的是庆长小判；每换一种小判，预测的米价就乘上庆长小判与新小判每两纯金之比。哪一年算作换了钱，要做一个假设：元禄小判从 1696 年起，乾字小判从 1711 年起；1714 年以后的好小判，要到元禄、乾字小判在 1718 年与 1720 年被禁用以后才成为主流，所以享保小判从 1719 年算起；元文小判从 1737 年算起。表~\ref{tab:toku-test} 与图~\ref{fig:toku_rice}(a) 是结果。

\begin{table}[htbp]\centering\small
\caption{按成色预测的米价与实际米价}\label{tab:toku-test}
\begin{tabularx}{\linewidth}{@{}l l Y Y Y@{}}
\toprule
& 时期（主要流通的钱） & 每单位钱的纯金或成色 & 预测 & 实际（平均） \\\midrule
\multicolumn{5}{@{}l}{江户（两/石）} \\
& 1685—1694（庆长小判） & 15.05 克 & 基准 & 0.786 \\
& 1696—1710（元禄小判） & 10.24 克 & 1.155 & 1.148 \\
& 1711—1718（乾字小判） & 7.91 克 & 1.496 & 1.781 \\
& 1719—1736（享保小判） & 15.49 克 & 0.763 & 0.838 \\
& 1737—1750（元文小判） & 8.63 克 & 1.371 & 1.154 \\\midrule
\multicolumn{5}{@{}l}{大阪（匁/石）} \\
& 1701—1705（元禄银） & 64\% & 基准 & 72.4 \\
& 1706—1710（宝永二宝银） & 50\% & 92.7 & 71.5 \\
& 1711—1718（四宝银） & 20\% & 231.8 & 118.7 \\
& 1719—1736（享保银） & 80\% & 58.0 & 47.1 \\
& 1737—1750（元文银） & 46\% & 100.8 & 68.0 \\
\bottomrule
\end{tabularx}
\datasource{米价据岩桥（1981）附表 1，经 GPIH 整理（\texttt{gpih\_iwahashi\_rice}）；成色据表~\ref{tab:toku-coins}。预测 = 基准 × 基准期的纯金（或成色）/ 该期的纯金（或成色），由本书计算。}
\end{table}

元禄小判的预测几乎与实际一样：1.155 对 1.148。乾字小判时期的实际米价高于预测，享保时期也略高。元文改铸是一个反例：按成色，米价应比享保时期涨 80\%（$15.49/8.63\approx1.795$），实际只涨了 38\%（$1.154/0.838\approx1.376$）；数量论的预测也偏高：假定旧币全部按铸币厂的 100 换 165 兑换，钱的数量增加 65\%，米价也应涨 65\%。可能的原因有三。一是新旧钱一起流通：Suzuki 说，1736 年市面上同时流通着六种小判与八种丁银\cite{suzuki2021}，法令又规定新旧等值，旧的好钱没有全部换成新钱，市面上钱的平均成色降得没有那么多。二是米的金属价格 $q$ 本身在变：享保年间连年丰收，米在实际上也便宜了\cite{schaede1989}，此后未必一样。三是江户的米价是官府规定的价格：按斋藤修（Saito Osamu）为数据写的说明，岩桥认为这类价格长期看接近市价，但年与年之间的波动“被压低了”\cite{iwahashi1981}。

\fig{toku_rice}{德川日本的米价与钱币的成色，1680—1760 年。(a) 看实线是否跟着褐色虚线的台阶走：实线是江户的米价（两/石），褐色虚线是按式~\eqref{eq:toku-content} 从小判的纯金算出的预测，以 1685—1694 年的平均为基准；元禄小判时期两者几乎重合，元文改铸以后实际米价远低于预测。(b) 看实线与褐色虚线的距离：实线是大阪堂岛的米价（匁/石），褐色虚线是按表~\ref{tab:toku-test} 各时期银币的成色算出的预测，以 1701—1705 年（元禄银，64\%）的平均为基准；预测在每一个时期都偏高，1711—1718 年偏高近一倍。纵轴为对数刻度。浅色带是 1711—1718 年钱币最差的时期。数据：岩桥（1981）附表 1，经 GPIH 整理；成色据 Suzuki（2021）表 2、表 3。}{fig:toku_rice}

\subsubsection{大阪的银价米}
大阪的银价米给成色说出了更大的难题（图~\ref{fig:toku_rice}(b)）。以 1701—1705 年的元禄银为基准，按表~\ref{tab:toku-test} 各时期银币的成色，四宝银时期的米价应为 231.8 匁，实际平均只有 118.7 匁；每一个时期的预测都偏高。最可能的原因，还是新旧钱一起流通：四宝银只铸了约 40 万贯，此前的元禄银、二宝银、三宝银加起来还有一百多万贯（表~\ref{tab:toku-coins}），市面上的银子并不都是 20\% 的四宝银。成色说在转换的时点上表现得最好：1718 年的法令规定好银 1 匁当四宝银 4 匁，米价从 1718 年的 127.9 匁降到 1719 年的 40.6 匁，降为 $1/3.15$，与 $1/4$ 相去不远\cite{iwahashi1981,suzuki2021}。

\subsubsection{清代：先涨的银价}
清代的数据检验的是式~\eqref{eq:toku-twoprice}。它的预测是：如果铜钱工资不变，银两工资就随银价下跌。图~\ref{fig:qing_cash}(b) 正是这样：1807—1813 年到 1845—1850 年，银价涨了 92\%，铜钱工资只涨了 4.8\%，银两工资跌了 45\%\cite{gamble1943}。它还给争论提供了一条线索：按 Gamble 的数据，1807—1824 年银价已经涨了 23\%，早于 Irigoin 所说的 20 年代中期才开始的白银外流（\secref{sec:toku-debate}）。这个检验有它的限度：数据只有北京郊区一处，工资是一种职业的工资；1853 年以后有了大钱，比价的单位变了，本章不用 1853 年以后的数字。

\begin{deeper}[史料与几处不一致]
\textbf{两种米价。}Schaede 引 Sugie（1984）的研究说，幕府提高币值以后，“1714 年卖 200 匁的一石米，1718 年只卖约 30 匁”\cite{schaede1989}；岩桥的大阪米价是 1714 年 149.2 匁、1718 年 127.9 匁、1719 年 40.6 匁\cite{iwahashi1981}。两者的年份与数值都对不上，本章用岩桥的年度数据。岩桥的大阪米价是堂岛的肥后米价格，江户米价则是幕府规定的价格\cite{iwahashi1981}。

\textbf{Suzuki 的几个数。}按她的表 2，庆长小判含金 84.291\%、重 17.85 克，元禄小判 57.366\%、17.85 克，纯金减少 31.94\%，她的正文写作 31.91\%；享保小判到元文小判，按表 2 是 44.30\%，正文写作 44.34\%\cite{suzuki2021}。越后屋一例里的纯金（新小判 8.581726 克，0.8 两旧小判 11.9759232 克）隐含的重量约为 13.06 克与 17.76 克，与表 2 的 13.13 克与 17.85 克略有出入。新井白石的改铸数字，金币每两旧币铸出 1.500 两新币，比按成色算的 1.469 两略多；银币 1.243，按成色是 1.25。铸币厂官员的总数意味着金币 44\%、银币 72\% 被改铸，Suzuki 概括为“大体上一半”。Suzuki 的表是图片，本书按页面图像读出。

\textbf{Schaede 的保证金。}他说 1770 年的保证金是每石“一到两匁”，又说是每张 100 石的合约 200 到 300 匁；按每张 100 石算，前者只有 100 到 200 匁，两个说法对不上。按 1770 年 66.3 匁的米价，300 匁约为 4.5\%，他写作“只有 5\%”。

\textbf{清代史料的日期与断句。}刘良驹的疏，吴嘉宾说是“道光乙巳年”（1845）\cite{xubiansheng}。曾国藩的两道奏疏，数字文本没有写明年月：《备陈民间疾苦疏》说江苏“自十六年至今。岁岁报歉”，写时他在刑部任职（“臣在刑部”），《拟银钱并用章程疏》提到道光二十四、二十五、二十六年的三道奏疏，可见写在道光二十六年（1846）以后\cite{xubiansheng}。刘良驹的一句，数字文本断作“几及二千文外。省则”，应为“几及二千文。外省则”。包世臣的《银荒小补说》，维基文库的《齐民四术》本第二处作“一千八百八十五文”，《续编》作“一千八百八十余文”。《清史稿》的“十年而罢”有两种读法（见案例）。乾隆年间京师 800 文的比价，《清史稿》写在乾隆十三年条下，以“先是”追述，只能说在 1748 年以前。

\textbf{Gamble 的工资。}GPIH 的工作簿把 Gamble 的工资加了倍，以计入雇主提供的伙食（按 Allen 等人的附录 II）\cite{allen2011china}；这不影响图~\ref{fig:qing_cash}(b) 的指数。比价在 1853 年以后跳到 2,723 文（1854）与 4,970 文（1856），那是大钱的年代，工作簿没有说明记账单位，本章不用。

\textbf{二手的转述。}林满红的书本章未能读到，她的观点据 Atwell 的书评与 Irigoin 的转述\cite{lin2006}；《管子》的话据 Suzuki 引的 Rickett 英译转译。
\end{deeper}

\section{今天还剩什么}\label{sec:toku-legacy}

堂岛米市与清代的银两早已不在了，这一章的有些想法也被推翻了：周腾虎以为比价的变化“何损于中国”，忽略了以名义单位固定下来的税与债（\secref{sec:toku-debate}）；“银贵只因白银外流”的单一解释，今天的研究也不接受。另外几个问题今天还在。

\textbf{期货、保证金与中央清算。}标准合约、保证金、定期结算与清算所担保，是今天期货市场的基本制度。保证金太低时清算所要承担的风险（例~\ref{ex:toku-hedge}），与全球金融危机以后把场外衍生品集中到中央对手方清算的讨论是同一个问题（\chref{ch:gfc}）。

\textbf{铸币税与贬值。}元禄改铸是通过降低成色征的一笔税，铸币税率约 32\%；今天的政府不再改铸金币，而是发行货币，通胀税的道理相同（\chref{ch:fiscaldominance}，\chref{ch:hyperinflation}）。

\textbf{货币错配。}清代的纳税人收入是铜钱，负担的却是银两；银贵就像本币贬值、外币债务加重。今天的新兴市场国家借外币、收本币，汇率贬值时也面临同样的“暗加”的负担（\chref{ch:emcrises}，\chref{ch:hierarchy}）。

\textbf{名义合约与通缩。}元文改铸用新币还旧债，减轻了大名与武士的债务；银贵让以银计的债务与税额变重。名义债务在通胀与通缩中的再分配，是费雪债务通缩理论的主题（\chref{ch:depressionus}）。

\textbf{菜单成本。}越后屋的 60 匁换算表，是为了避免频繁改价；改价的成本为什么会让价格变得粘滞，是\chref{ch:newkeynesian}的问题。

\textbf{记账单位。}米在德川日本不是交换媒介，却是基本的记账单位。Schaede 说，要完整把握米的作用，也许需要一个新的货币概念，它允许一种东西“数量既不受限也不可控，却仍是经济与经济政策的基本计量单位”\cite{schaede1989}。货币的三种职能可以分开，见\chref{ch:money}。

\begin{takeaways}
\begin{itemize}
\item 德川日本按米计量土地与赋税，却用三种钱：东边的金币按面值计数，西边的银子按重量计价，铜钱按枚。同一件货的两种价格满足 $P_a=e\,P_b$，式~\eqref{eq:toku-twoprice}；1700 年以后江户按 60 匁折算，西边的比价是 45 匁时，江户的标价就高 33\%。以米计税，米价的风险在收税的一方，所以“米将军”想要米贵；以银计税，风险在纳税的一方，$x=Te/p$，式~\eqref{eq:toku-burden}。
\item 元禄改铸（1695）把小判的纯金从 15.05 克降到 10.24 克。按成色说，$P=q/c$，式~\eqref{eq:toku-content}，江户的米价应涨 46.9\%，实际 1696—1710 年比 1685—1694 年涨了 46\%；新井白石估计流通量没有变化，数量论解释不了这次通胀。幕府的铸币税，式~\eqref{eq:toku-seign}，约 432 万两，税率约 32\%；元文改铸（1736）的铸币厂按 100 两旧币换 165 两新币，税率只有 5\%—8\%。
\item 堂岛米市在 1730 年得到承认：米票从仓库的提单变成了标准化的期货合约，100 石一张，清算所收保证金、担保履约，十天结算一次。卖出期货可以锁定 $QF_0$，式~\eqref{eq:toku-hedge}；保证金从三分之一降到 5\%，同样的价格变动对保证金的冲击就从 30\% 变成 200\%。1798—1835 年的检验不能拒绝期货价格是到期现货价格的无偏预测。
\item 清代的银钱比价从 18 世纪的七八百文、一千文上下，涨到 1845—1850 年平均约 1,990 文。以钱计的米价与工资几乎不动，以银计的就跌了一半：曾国藩说“小民暗加一倍之赋”，北京郊区的银两工资跌了 45\%。
\item 银贵的原因当时就有争论：白银外流（黄爵滋、林则徐）、窖藏与“用银太重”（刘良驹、吴嘉宾），还是银贵无害（周腾虎）。今天的研究强调拉丁美洲白银供给的中断、世界需求的下降（林满红）与好银元的短缺（Irigoin）；银价在外流开始以前就在上涨。
\end{itemize}
\end{takeaways}

\begin{selfcheck}
\begin{enumerate}
\item 一种金币每两的纯金从 15 克降到 12 克。按式~\eqref{eq:toku-content}，以两计的物价会涨多少？如果铸币厂每收回 100 两旧币付给 105 两新币，按式~\eqref{eq:toku-seign}，每 100 两旧币的铸币税是多少两？税率是多少？
\item 1700 年以后，江户按 1 两 60 匁折算。如果大阪的比价是 1 两 50 匁，一件在大阪值 100 匁的货，在江户卖多少两？标价多少匁？比大阪贵百分之几？
\item 一个纳税人每亩的税是 0.5 两银子，米价每石 2,500 文。银钱比价从 1,200 文涨到 1,800 文，按式~\eqref{eq:toku-burden}，他每亩要卖的米从多少变成多少？如果米价同时跌了 10\% 呢？
\item 一个米商持有 200 石米的米票，按 50 匁卖出两张期货合约。到期时现货 45 匁，期货 45.5 匁。按式~\eqref{eq:toku-hedge}，他的米票与期货合计值多少？没有对冲时少了多少？
\item 一个工人的铜钱工资每天 90 文，银钱比价从 1,100 文涨到 2,000 文。他的工资折成银两跌了多少？如果他买的东西都以铜钱计价，价格不变，他的生活变坏了吗？什么情况下会变坏？
\item 设人口 3.5 亿，吸食鸦片的人占 0.5\%，每人每天花银一钱，一年按 360 天。按林则徐的算法，一年花多少银子？这能说明一年有这么多白银流出中国吗？
\end{enumerate}
\end{selfcheck}

\begin{answers}
\item 物价乘上 $15/12=1.25$，涨 25\%。新币 $100\times15/12=125$ 两，$S=100\times(1.25-1-0.05)=20$ 两；税率 $1-1.05/1.25=16\%$，即 20/125。
\item $100/50=2$ 两；江户按 60 匁折算，标价 120 匁，比大阪贵 20\%，即 $\bar e/e_W=60/50$。
\item $x=0.5\times1200/2500=0.24$ 石，变成 $0.5\times1800/2500=0.36$ 石，多了一半；米价再跌 10\%，$0.5\times1800/2250=0.4$ 石。
\item $V_T=200\times50+200\times(45-45.5)=9{,}900$ 匁；没有对冲时只值 $200\times45=9{,}000$ 匁，少了 900 匁。对冲以后，只剩下 0.5 匁的基差造成的 100 匁损失。
\item $90/1100\approx0.082$ 两降到 $90/2000=0.045$ 两，跌了 $1-1100/2000=45\%$。买的东西都以铜钱计价、价格不变，他的生活没有变坏；要用银子付的东西，比如按银两计的税或租，才让他变坏。
\item $3.5\times10^8\times0.5\%\times36=6{,}300$ 万两。不能：这是花在鸦片上的钱，不是流出的白银；鸦片的价钱里有国内的转运与零售，进口也可以用出口货物抵付，只有贸易差额需要用白银结清。
\end{answers}

斯密写《国富论》的 1776 年，堂岛米市已经开了 46 年，清朝的钱价按吴嘉宾的说法还在一千一百文以内。这一章的市场都在没有经济学理论的情况下运转；下一章去看第一个系统地追问国民财富从哪里来、银行应当按什么原则放贷的人：亚当·斯密。
